Anna et Baptiste jouent au jeu de l'\emph{inékoalaté}, un jeu à deux
joueurs dont les règles dépendent d'un réel $\lambda > 0$ connu des
deux joueurs. Pour tout $n \geqslant 1$, lors du $n^{\text{ème}}$ tour
de jeu,
\begin{itemize}
  \item si $n$ est impair, Anna choisit un réel $x_n \geqslant 0$ tel
  que
  \[
  x_1 + x_2 + \cdots + x_n \leqslant \lambda n ;
  \]
  \item si $n$ est pair, Baptiste choisit un réel $x_n \geqslant 0$
  tel que
  \[
  x_1^2 + x_2^2 + \cdots + x_n^2 \leqslant n.
  \]
\end{itemize}
Si un joueur ne peut pas choisir de réel $x_n$ adéquat, la partie
s'arrête ; il perd et son adversaire gagne. Si la partie ne s'arrête
jamais, aucun joueur ne gagne. À tout moment, chaque joueur connaît la
liste de tous les réels déjà choisis.

Trouver toutes les valeurs de $\lambda$ pour lesquelles Anna dispose
d'une stratégie gagnante, ainsi que toutes celles pour lesquelles
Baptiste dispose d'une stratégie gagnante.
